RISC-V is gaining popularity as a free and open-source ISA. PicoRV32 is a small RISC-V soft core CPU designed to be deployed on FPGAs with the fewest hardware resource requirements. However, the arithmetic units of the core focus on hardware size requirements and can be a bottleneck when it comes to processing speed. This is a problem when implementing edge-AI style inference processing, as the primary operation is real-time processing dominated by loops of multiply-accumulate instructions with tight power and latency requirements. This paper proposes the design of an arithmetic unit core optimization using a 32-bit Kogge-Stone adder design and a Vedic multiplier using the Urdhva-Tiryakbhyam algorithm. The multiplier is integrated with the Pico Co-Processor Interface (PCPI) as a three-stage pipeline structure. Experimental results using the project dataset show a reduction in multiplication latency from 36 cycles to 3 cycles at 100~MHz operation, giving a 12.0$\times$ speedup factor. Additionally, the logic depth of the ALU add operation is reduced from 32 to 6 stages. Formal verification results show 0 failures from 70,598 simulation vectors with 13/13 GF(2) checks successful, validating the optimization design functionality. These results place the optimized core as a viable compute IP option for edge-AI SoCs with a need for efficient and verifiable integer arithmetic.
